% TOP_PROBLEM: {"id": 30006895, "title": "Toms–Winter: Strict Comparison Implies Z-Stability", "category_id": 8, "category": {"id": 8, "name": "analysis", "display_name": "Analysis", "description": "Limits, continuity, calculus, and function theory.", "slug": "analysis", "order_index": 8, "created_at": "2026-07-31T15:26:25.671Z"}, "status": "open"}
% FIRST_POSTED: 2026-09-25
% ENTRY_KIND: statement_only
% !TeX program = lualatex
% Definition and sources only; this file is not an LLM solution attempt.
\documentclass[11pt]{article}
\usepackage[margin=25mm]{geometry}
\usepackage{fontspec}
\setmainfont{Latin Modern Roman}
\usepackage{amsmath,amssymb,mathtools}
\usepackage{unicode-math}
\setmathfont{Latin Modern Math}
\usepackage{xurl,hyperref}
\hypersetup{colorlinks=true,urlcolor=blue}
\setcounter{secnumdepth}{0}
\setlength{\parindent}{0pt}
\setlength{\parskip}{5pt}
\setlength{\emergencystretch}{3em}
\begin{document}
\section*{\texorpdfstring{Toms–Winter: Strict Comparison Implies Z-Stability}{Toms–Winter: Strict Comparison Implies Z-Stability}}\label{30006895}
\subsection{Definitions and mathematical statement}
Let $A$ be a separable simple nuclear C*-algebra, not isomorphic to an algebra of compact operators; no unit is required. Its Cuntz semigroup $\operatorname{Cu}(A)$ consists of positive elements of $A\otimes\mathcal K$ modulo mutual Cuntz subequivalence, where $a\precsim b$ means $x_j^*bx_j\to a$ for some sequence. In the selected strict-comparison convention, $x\le y$ whenever every additive order-preserving functional $\varphi:\operatorname{Cu}(A)\to[0,\infty]$, preserving zero and increasing sequential suprema, satisfies $\varphi(x)\le\varphi(y)$, with strict inequality whenever $0<\varphi(y)<\infty$. The conjecture asserts
\[
 A\text{ has strict comparison}\quad\Longrightarrow\quad
 A\cong A\otimes_{\min}\mathcal Z,
\]
where $\mathcal Z$ is the Jiang--Su algebra. Separability means a countable norm-dense subset; simplicity means no nontrivial closed two-sided ideals; nuclearity is the completely positive matrix-factorization approximation property. No UCT, finite nuclear dimension, or extra trace-space hypothesis is appended.
\par\smallskip\textit{Formulation and concept references:} \href{https://arxiv.org/pdf/1912.04207}{[S1]}.

\subsection{Short English statement}
Does strict comparison force every separable simple nuclear non-elementary C*-algebra to absorb the Jiang--Su algebra, including nonunital cases?
\subsection{Sources}
\begin{itemize}
\item[S1]Castillejos, Evington, Tikuisis and White, Uniform Property Gamma, Conjecture 5.1 and following comparison definition; Schafhauser, Tikuisis and White, Nuclear C*-algebras: 99 problems, Problem XVIII.
\url{https://arxiv.org/pdf/1912.04207}.
\textit{Formulation source.}
\end{itemize}
\end{document}
